\begin{remark}[Basic-vector coefficients and nearest-neighbor local-term commutation]%
    \label{rem:product_pair_projectors}
    \lean{MPSTensor.productPairWindow}
    \lean{MPSTensor.productPairWindow_one}
    \lean{MPSTensor.productPairState}
    \lean{MPSTensor.productPairState_one}
    \lean{MPSTensor.HasProductPairMPV}
    \lean{MPSTensor.HasProductPairLocalProjectors}
    \lean{MPSTensor.HasProductPairLocalProjectors.of_commuting_localTerms}
    \lean{MPSTensor.HasProductPairLocalProjectors.isNNCPH}
    \lean{MPSTensor.ProductPairBridge}
    \lean{MPSTensor.ProductPairBridge.localTerm_idempotent}
    \lean{MPSTensor.ProductPairBridge.commuting_twoSite_localTerms}
    \lean{MPSTensor.ProductPairBridge.isNNCPH}
    \lean{MPSTensor.AppendixBStructuralData}
    \lean{MPSTensor.AppendixBStructuralData.exists_ofRFP}
    \lean{MPSTensor.AppendixBStructuralData.ofRFP}
    \lean{MPSTensor.AppendixBStructuralData.twoSiteAmplitude}
    \lean{MPSTensor.AppendixBStructuralData.coreTensor}
    \lean{MPSTensor.AppendixBStructuralData.coreTensor_apply}
    \lean{MPSTensor.AppendixBStructuralData.coreTensor_pair_orthogonality}
    \lean{MPSTensor.cyclicVirtualSucc}
    \lean{MPSTensor.AppendixBStructuralData.cyclicVirtualPairState}
    \lean{MPSTensor.AppendixBStructuralData.mpv_coreTensor_eq_cyclicVirtualPairState}
    \lean{MPSTensor.AppendixBStructuralData.cyclicVirtualPairState_two_eq_twoSiteAmplitude}
    \lean{MPSTensor.trace_gauge_boundary}
    \lean{MPSTensor.GaugeEquiv.groundSpace_le}
    \lean{MPSTensor.GaugeEquiv.groundSpace_eq}
    \lean{MPSTensor.AppendixBStructuralData.gaugeEquiv_coreTensor}
    \lean{MPSTensor.AppendixBStructuralData.groundSpace_eq_coreTensor}
    \lean{MPSTensor.AppendixBStructuralData.mpv_eq_coreTensor}
    \lean{MPSTensor.AppendixBStructuralData.mpv_eq_cyclicVirtualPairState}
    \lean{MPSTensor.AppendixBStructuralData.twoSiteAmplitude_eq_mpv}
    \lean{MPSTensor.AppendixBStructuralData.twoSiteAmplitude_eq_coreTensor_mpv}
    \lean{MPSTensor.AppendixBStructuralData.mpv_coreTensor_eq_productPairState_one}
    \lean{MPSTensor.AppendixBStructuralData.mpv_eq_productPairState_one}
    \lean{MPSTensor.AppendixBProductPairExtraction}
    \lean{MPSTensor.AppendixBProductPairExtraction.ofCoreTensorFactorization}
    \lean{MPSTensor.AppendixBProductPairExtraction.ofCoreTensorFactorizationAndCommutation}
    \lean{MPSTensor.AppendixBProductPairExtraction.ofCoreTensorFactorizationAndOverlapCommutation}
    \lean{MPSTensor.AppendixBProductPairExtraction.ofCoreTensorFactorizationAndAdjacentCommutation}
    \lean{MPSTensor.AppendixBProductPairExtraction.toProductPairBridge}
    \lean{MPSTensor.AppendixBProductPairExtraction.commuting_twoSite_localTerms}
    \lean{MPSTensor.commuting_twoSite_localTerms_of_rfp_of_appendixBExtraction}
    \lean{MPSTensor.rfp_implies_nncph_of_appendixBExtraction}
    \lean{MPSTensor.rfp_implies_nncph_ground_state_of_appendixBExtraction}
    \leanok
    \uses{thm:trace_eval_word_eq_sum_cyclic,
        lem:appendixB_basis_change_parent_interaction,
        def:appendixB_two_site_basic_space, def:appendixB_qax,
        def:appendixB_qxb, lem:appendixB_qax_qxb_parent_interaction,
        lem:appendixB_three_site_parent_lifts,
        lem:appendixB_lift_commutator_parent_terms, lem:local_term_idempotent}
    The Appendix~B structural form supplies $A^i=X\Lambda U^iX^{-1}$.
    In the source unit pair-index convention, the residual tensor satisfies
    \begin{align}
        \sum_i
        \overline{(U^i)_{\alpha,\beta}}
        (U^i)_{\alpha',\beta'}
        &=
        \delta_{\alpha,\alpha'}\delta_{\beta,\beta'}.
        \notag
    \end{align}
    Therefore the residual core tensor $\Lambda U^i$ satisfies the weighted
    pair-index relation
    \begin{align}
        \sum_i
        \overline{(\Lambda U^i)_{\alpha,\beta}}
        (\Lambda U^i)_{\alpha',\beta'}
        &=
        \lambda_\alpha\lambda_{\alpha'}
        \delta_{\alpha,\alpha'}\delta_{\beta,\beta'}.
        \notag
    \end{align}
    This is an orthogonality relation after summing over the physical letter
    $i$. It is not a coefficient identity for a fixed physical word
    $\sigma$.
    The coefficient associated to the source basic-vector formula is
    \begin{align}
        C_L(\sigma)
        &=
        \sum_{\alpha_0,\ldots,\alpha_{L-1}}
        \prod_{t=0}^{L-1}
        \Lambda_{\alpha_t}
        (U^{\sigma_t})_{\alpha_t,\alpha_{t+1}},
        \qquad \alpha_L=\alpha_0.
        \notag
    \end{align}
    This is the coefficient expression for the source vector
    \begin{align}
        \ket{V^{(L)}(A_j)}
        &=
        U^{\otimes L}\ket{\varphi_j}^{\otimes L},
        \notag\\
        \ket{\varphi_j}
        &=
        \sum_m\lambda_m\ket{m,m},
        \notag
    \end{align}
    where $\varphi_j$ is shared by $b_t$ and $a_{t+1}$ and $U$ acts on
    $(a_t,b_t)$.
    At length two, this cyclic coefficient is the structural two-site amplitude:
    $C_2(\sigma)=\phi_2(\sigma)$.
    The change of basis does not alter either the local spaces $G_L$ or the
    periodic MPS coefficients: $G_L(A)=G_L(\Lambda U)$.
    It therefore also leaves the canonical parent interaction unchanged:
    $q_L(A)=q_L(\Lambda U)$.
    At length two, the Appendix~B two-site basic support is
    $G_2(\Lambda U)$, and the coefficient-space representatives
    $\widehat Q_{AX}$ and $\widehat Q_{XB}$ are the two copies of
    $q_2(\Lambda U)$ before they are lifted to the $AX$ and $XB$ faces.
    Hence
    \begin{align}
        \widehat Q_{AX}
        &=
        q_2(A),
        \notag\\
        \widehat Q_{XB}
        &=
        q_2(A),
        \notag
    \end{align}
    as coefficient-space representatives. On the three-site window their two
    lifts give the local identifications
    \begin{align}
        h_0^{(3)}(A,2)
        &=
        (\widehat Q_{AX})_{AX},
        \notag\\
        h_1^{(3)}(A,2)
        &=
        (\widehat Q_{XB})_{XB}.
        \notag
    \end{align}
    Consequently, for every chain length $L\geq1$, the original tensor has
    the same cyclic coefficient formula:
    $V^{(L)}(A)_\sigma=C_L(\sigma)$.
    In particular, the corresponding two-site amplitude is
    \begin{align}
        \phi_2(a,b)
        &=
        V^{(2)}(A)(a,b) =
        V^{(2)}(\Lambda U)(a,b).
        \notag
    \end{align}
    Thus the one-pair case of the factorization is already verified for the core
    tensor:
    $V^{(2)}(\Lambda U)(\sigma)=\phi_2(\sigma_0,\sigma_1)$.
    The even-chain physical-pair factorization is the following condition, for a
    positive number $n$ of pairs,
    \begin{align}
        V^{(2n)}(A)(\sigma)
        &=
        \prod_{p=0}^{n-1}
        \phi_2(\sigma_{2p},\sigma_{2p+1}).
        \notag
    \end{align}
    The case $n=0$ is not part of this condition: the empty-chain MPV
    coefficient is the virtual bond dimension, whereas the empty product on the
    right-hand side is $1$.
    The first equality, $V^{(L)}(\Lambda U)_\sigma=C_L(\sigma)$, is the trace
    expansion of the closed matrix product and is the coefficient form of
    $U^{\otimes L}\varphi_j^{\otimes L}$. The remaining coefficient step is
    to relate this source expression, with the physical word $\sigma$ fixed, to
    the even-chain physical-pair factorization above.
    On a three-site chain, the two adjacent length-two local terms are the
    $AX$ and $XB$ actions of the canonical two-site parent interaction.
    Theorem~\ref{thm:appendixB_overlapping_two_site_commutation} constructs the
    adjacent physical support projectors from the basic-vector expression and
    proves the overlapping relation
    \begin{align}
        (\widehat Q_{AX})_{AX}(\widehat Q_{XB})_{XB}
        &=
        (\widehat Q_{XB})_{XB}(\widehat Q_{AX})_{AX}.
        \notag
    \end{align}
    Theorem~\ref{thm:three_site_commutator_chain_transport} transports this
    identity around a periodic chain. Consequently, for every $N>2$ and every
    pair of sites $i,j\in\Z/N\Z$,
    $h_i(A,2)h_j(A,2) = h_j(A,2)h_i(A,2)$.
    Definition~\ref{def:appendixB_chain_local_projectors} collects these maps
    into a commuting idempotent family; idempotency is
    Lemma~\ref{lem:local_term_idempotent}. Together with
    Lemma~\ref{lem:parent_hamiltonian_ff}, this gives the zero-energy equation
    for $V^{(N)}(A)$. Neither conclusion requires the even-chain
    physical-pair factorization displayed above. That factorization remains a
    separate conditional coefficient statement. The full source implication
    additionally requires removal of the left-canonical hypothesis and the
    ground-space spanning equation. The coefficient-space analog of the
    local kernel-intersection equation is
    Theorem~\ref{thm:appendixB_definition_d2_coefficient_analogue}; identifying
    its hatted maps with the source projectors remains separate and is not
    needed for the chain commutator.
\end{remark}

\begin{theorem}[Even-chain factorization and ground-space spanning give all-chain NNCPH]%
    \label{thm:appendixB_product_pair_has_nncph_ground_spaces}
    \lean{MPSTensor.ProductPairBridge.hasNNCPHGroundSpaces_of_groundSpaceSpanning}
    \lean{MPSTensor.rfp_implies_hasNNCPHGroundSpaces_of_appendixBExtraction_of_groundSpaceSpanning}
    \leanok
    \uses{rem:product_pair_projectors, def:has_nncph_ground_spaces,
        def:parent_hamiltonian_ground_space_spanning}
    Suppose the conditional even-chain hypotheses give commuting two-site parent
    terms for a tensor $B$. If the nearest-neighbor parent Hamiltonian also
    satisfies the ground-space spanning equation of
    Definition~\ref{def:parent_hamiltonian_ground_space_spanning}, with BNT
    components $A_1,\ldots,A_g$, then $B$ satisfies the all-chain
    nearest-neighbor commuting parent-Hamiltonian ground-space condition of
    Definition~\ref{def:has_nncph_ground_spaces}.

    In particular, if a normal left-canonical RFP tensor has nonzero bond
    dimension, the same conditional Appendix~B hypotheses, and the same
    ground-space spanning equation, then it satisfies the full all-chain condition in
    \cite[Theorem~3.10(iii)]{Cirac2016MPDO_arXiv}.
\end{theorem}

\begin{proof}
    \uses{lem:parent_hamiltonian_ff, def:has_nncph_ground_space}
    \leanok
    The conditional hypotheses give $h_i h_j=h_j h_i$ for all translated
    two-site local terms. Lemma~\ref{lem:parent_hamiltonian_ff} gives
    $h_iV^{(N)}(B)=0$. For every $N>2$, the spanning hypothesis supplies
    $\ker H_2^{(N)}(B) = \spn\{V^{(N)}(A_j):j=1,\ldots,g\}$.
    These are exactly the three clauses in
    Definition~\ref{def:has_nncph_ground_space}, uniformly in $N$.
\end{proof}

\begin{theorem}[Eventual dimension of nearest-neighbor parent ground spaces]%
    \label{thm:nncph_eventual_ground_space_dimension}
    \lean{MPSTensor.HasBNTSectorData.eventually_parentHamiltonian_groundSpace_finrank_eq}
    \leanok
    \uses{def:sector_weight_data, def:has_nncph_ground_spaces}
    Let $P$ be a sector decomposition with representatives
    $A_1,\ldots,A_g$ and total tensor $B$. Suppose that the vectors
    $\{\ket{V^{(N)}(A_j)}\}_{j=1}^g$ are linearly independent for all
    sufficiently large $N$, and that $B$ satisfies the all-chain
    nearest-neighbor commuting parent-Hamiltonian ground-space condition with
    these representatives. Then there is $N_0$ such that, for every $N>N_0$,
    $\dim_{\C}\ker H_2^{(N)}(B) = g$.
    Thus the ground-state degeneracy is eventually constant, as required in
    \cite[Section~IV]{Beigi2012Classification}.
\end{theorem}

\begin{proof}\leanok
    \uses{def:bnt_span, def:has_nncph_ground_spaces}
    Choose $N_0\geq2$ beyond the linear-independence threshold. For $N>N_0$,
    the ground-space spanning equation gives
    $\ker H_2^{(N)}(B) = \spn\{\ket{V^{(N)}(A_j)}:j=1,\ldots,g\}$.
    Identifying these vectors with their coefficient functions preserves
    linear independence. The dimension of their span is therefore $g$.
\end{proof}

\begin{definition}[Two-site parent-interaction matrix]%
    \label{def:two_site_parent_interaction_matrix}
    \lean{MPSTensor.twoSiteParentInteractionMatrix}
    \leanok
    Let $A$ be a matrix product tensor. Write its canonical two-site parent
    interaction in the ordered-pair basis of the two-site physical Hilbert
    space, and denote the resulting matrix by $Q_A$.
\end{definition}

\begin{theorem}[Hermitian commuting overlapping parent interactions]%
    \label{thm:two_site_parent_interaction_commuting_overlaps}
    \lean{MPSTensor.IsNNCPH.twoSiteParentInteractionMatrix_isHermitian_and_overlappingLifts_commute}
    \leanok
    \uses{def:two_site_parent_interaction_matrix}
    Suppose that the canonical length-two parent terms of $A$ commute on a
    three-site chain. Then $Q_A$ is Hermitian, and its natural placements on
    the first and last pairs of a three-fold tensor product commute:
    \begin{align}
        (Q_A\otimes\Id)(\Id\otimes Q_A)
        &=
        (\Id\otimes Q_A)(Q_A\otimes\Id).
        \notag
    \end{align}
\end{theorem}

\begin{proof}\leanok
    The parent interaction is the orthogonal projector onto the orthogonal
    complement of the local matrix product space, hence is Hermitian.
    Transporting the adjacent three-site parent-term commutator to the common
    left-associated tensor-product basis gives the displayed identity.
\end{proof}

\begin{theorem}[Spatial block actions of the two-site parent interaction]%
    \label{thm:two_site_parent_interaction_spatial_blocks}
    \lean{MPSTensor.IsNNCPH.exists_unitary_blockActions_of_twoSiteParentInteractionMatrix}
    \leanok
    \uses{thm:two_site_parent_interaction_commuting_overlaps}
    Under the preceding hypothesis, there are positive integers
    $d_q,m_q$, a unitary identification
    \begin{align}
        \mathcal H
        &\cong
        \bigoplus_q
        \mathcal H_{q,l}\otimes\mathcal H_{q,r},
        \notag
    \end{align}
    and Hermitian matrices $R_q$ on
    $\mathcal H\otimes\mathcal H_{q,l}$ and $S_q$ on
    $\mathcal H_{q,r}\otimes\mathcal H$ such that, after conjugating the
    middle site by this unitary,
    \begin{align}
        Q_A
        &=
        \bigoplus_q(R_q\otimes\Id_{q,r})
        \qquad\text{on the first pair},
        \notag\\
        Q_A
        &=
        \bigoplus_q(\Id_{q,l}\otimes S_q)
        \qquad\text{on the last pair}.
        \notag
    \end{align}
    This is the spatial decomposition of
    \cite[Lemma~2.1]{Beigi2012Classification} for the canonical parent
    interaction.
\end{theorem}

\begin{proof}\leanok
    Apply the finite-dimensional spatial decomposition for two Hermitian
    overlapping operators to $Q_A$ in both positions, using
    Theorem~\ref{thm:two_site_parent_interaction_commuting_overlaps}.
\end{proof}

\begin{definition}[Finite sector graph]\label{def:beigi_sector_graph}
    \lean{MPSTensor.BeigiSectorGraphData}
    \lean{MPSTensor.BeigiSectorGraphData.unitary_mul_conjTranspose}
    \lean{MPSTensor.BeigiSectorGraphData.edgeGroundSpace}
    \lean{MPSTensor.BeigiSectorGraphData.IsEdge}
    \lean{MPSTensor.BeigiSectorGraphData.OrderedCycle}
    \leanok
    Suppose that a one-site unitary identifies the physical Hilbert space with
    a finite direct sum
    \begin{align}
        \mathcal H
        &\cong
        \bigoplus_a
        \mathcal H_{a,l}\otimes\mathcal H_{a,r},
        \notag
    \end{align}
    and that in these coordinates the complementary two-site interaction is
    $\bigoplus_{a,b}\Id_{a,l}\otimes P_{a,b}\otimes\Id_{b,r}$, where each
    $P_{a,b}$ is an orthogonal projector. The directed edge
    $a\to b$ is present precisely when $\ran P_{a,b}$ is nonzero.
    An ordered cycle of length $N$ is a sequence $(a_0,\ldots,a_{N-1})$ for
    which every cyclically adjacent edge $a_n\to a_{n+1}$ is present. Cyclic
    rotations are not identified. This is the graph following equation~(2)
    of \cite[Section~III]{Beigi2012Classification}.
\end{definition}

\begin{theorem}[Symmetric spatial form of a commuting positive interaction]
    \label{thm:commuting_positive_interaction_symmetric_spatial_form}
    \lean{Matrix.exists_positive_etaPair_decomposition_of_overlappingLifts_commute}
    \lean{Matrix.etaPairSpatialBlockEquiv}
    \leanok
    Let $B\geq0$ be a two-site operator whose overlapping placements commute:
    $(B\otimes\Id)(\Id\otimes B) = (\Id\otimes B)(B\otimes\Id)$.
    There are positive integers $d_{q,l},d_{q,r}$, a one-site unitary $U$,
    and positive semidefinite operators $\eta_{q,h}$ on
    $\mathcal H_{q,r}\otimes\mathcal H_{h,l}$ such that
    \begin{align}
        (U\otimes U)^*B(U\otimes U)
        &=
        \bigoplus_{q,h}
        \Id_{q,l}\otimes\eta_{q,h}\otimes\Id_{h,r}.
        \notag
    \end{align}
    This is the symmetric two-sided form obtained from
    \cite[Lemma~2.1 and Section~III, equation~(2)]{Beigi2012Classification}.
\end{theorem}

\begin{proof}\leanok
    Apply the spatial decomposition to the two overlapping placements of $B$.
    In the resulting coordinates, matrix entries vanish unless both the left
    and right sector labels agree. For fixed sectors $q,h$, choose unit
    vectors $l_0\in\mathcal H_{q,l}$ and
    $r_0\in\mathcal H_{h,r}$ and define
    \begin{align}
        \bra{r,l}\eta_{q,h}\ket{r',l'}
        &=
        \bra{l_0,r;l,r_0}\widetilde B\ket{l_0,r';l',r_0},
        \notag\\
        \widetilde B
        &=
        (U\otimes U)^*B(U\otimes U).
        \notag
    \end{align}
    The left block action makes the right outer factor an identity, while the
    right block action makes the left outer factor an identity. Hence the
    displayed matrix element is independent of the chosen outer coordinates
    and gives the stated direct sum. Each $\eta_{q,h}$ is a compression of
    the positive semidefinite matrix $\widetilde B$, so it is positive
    semidefinite.
\end{proof}

\begin{theorem}[Sector graph of a commuting parent interaction]
    \label{thm:nncph_beigi_sector_graph}
    \lean{MPSTensor.IsNNCPH.nonempty_beigiSectorGraphData}
    \lean{MPSTensor.IsNNCPH.beigiSectorGraphData}
    \lean{MPSTensor.twoSiteParentInteractionMatrix_isStarProjection}
    \lean{MPSTensor.twoSiteParentGroundProjectorMatrix_isStarProjection}
    \lean{MPSTensor.twoSiteParentGroundProjectorMatrix_posSemidef}
    \lean{MPSTensor.IsNNCPH.twoSiteParentGroundProjectorMatrix_overlappingLifts_commute}
    \leanok
    \uses{def:beigi_sector_graph}
    Suppose that the canonical length-two parent terms of $A$ commute on a
    three-site chain. Then the complementary two-site interaction admits a
    finite sector graph as in Definition~\ref{def:beigi_sector_graph}.
\end{theorem}

\begin{proof}\leanok
    \uses{thm:commuting_positive_interaction_symmetric_spatial_form}
    The canonical two-site parent interaction $Q_A$ is an orthogonal
    projector. Thus $\Id-Q_A$ is a positive semidefinite orthogonal
    projector. The overlapping placements of $\Id-Q_A$ commute because
    those of $Q_A$ commute. Apply
    Theorem~\ref{thm:commuting_positive_interaction_symmetric_spatial_form} to
    obtain
    \begin{align}
        (U\otimes U)^*(\Id-Q_A)(U\otimes U)
        &=
        \bigoplus_{q,h}
        \Id_{q,l}\otimes\eta_{q,h}\otimes\Id_{h,r}.
        \notag
    \end{align}
    The left-hand side is idempotent, so every block on the right is
    idempotent. Evaluating the two outer identity factors at fixed unit
    vectors gives $\eta_{q,h}^2=\eta_{q,h}$. Since $\eta_{q,h}\geq0$, each
    $\eta_{q,h}$ is an orthogonal projector and therefore defines the required
    edge projection.
\end{proof}

\begin{theorem}[Ordered-cycle ground-space dimension formula]
    \label{thm:beigi_ordered_cycle_ground_space_dimension}
    \lean{MPSTensor.BeigiSectorGraphData.parentHamiltonianGroundSpaceES_finrank_eq}
    \leanok
    \uses{def:beigi_sector_graph}
    For every $N\geq2$, the periodic ground-space dimension is
    \begin{align}
        \dim\ker H_N
        &=
        \sum_{(a_0,\ldots,a_{N-1})}
        \prod_{n=0}^{N-1}
        \dim\ran P_{a_n,a_{n+1}},
        \notag
    \end{align}
    where the sum is over ordered cycles and the indices are read modulo $N$.
    This is equations~(3)--(4) of
    \cite[Section~III]{Beigi2012Classification} for chains of length at least
    two. The one-site periodic interaction requires a separate convention.
\end{theorem}

\begin{proof}\leanok
    The translated parent interactions commute. Hence the ordered product of
    their complementary orthogonal projectors is the projector onto
    $\ker H_N$. In the cyclic sector coordinates this product is block
    diagonal, with the block indexed by $(a_0,\ldots,a_{N-1})$ equal to the
    tensor product of the adjacent projectors $P_{a_n,a_{n+1}}$. Taking the
    trace gives the product of their ranks. A sequence containing a missing
    edge contributes zero, leaving exactly the displayed ordered-cycle sum.
\end{proof}

\begin{theorem}[Period-two ordered-cycle ground-space dimension formula]
    \label{thm:beigi_ordered_cycle_ground_space_dimension_two}
    \lean{MPSTensor.BeigiSectorGraphData.parentHamiltonianGroundSpaceES_finrank_eq_two}
    \leanok
    \uses{def:beigi_sector_graph}
    At period two, the ground-space dimension is
    \begin{align}
        \dim\ker H_2
        &=
        \sum_{\substack{a\to b\\ b\to a}}
        \dim\ran P_{a,b}
        \dim\ran P_{b,a}.
        \notag
    \end{align}
    The sum is over ordered pairs: if $a\ne b$, then $(a,b)$ and $(b,a)$
    are distinct summands. This is the length-two convention stated after
    equation~(4) of \cite[Section~III]{Beigi2012Classification}.
\end{theorem}

\begin{proof}\leanok
    \uses{thm:beigi_ordered_cycle_ground_space_dimension}
    Specialize the ordered-cycle formula to $N=2$. Its two cyclic factors are
    $P_{a,b}$ and $P_{b,a}$, which gives the displayed product of ranks.
\end{proof}
